\section{Proofs of Main Results}

In this chapter, we provide the proofs of the main results from Chapters~\ref{sec:proposed_algorithms} and \ref{sect_regret}.

\subsection{Proofs of Lemmas on Number of Resampling}
%\section{Proofs of proposed algorithms}
In this section, we provide the proofs of Lemmas~\ref{lem: general idea}--\ref{lem: complexity of CGR2}. 

For convenience of describing the conditions in the proof, we define an indicator function as
\begin{equation*}
    \chi_{t,i}(r_t') =
    \begin{cases} 
        1, & \text{if } i = \arg\min_{j \in [K]} \{\hat{L}_t - r_t' / \eta_t\}, \\
        0, & \text{otherwise}.
    \end{cases}
\end{equation*}
Then, the condition \eqref{cond_termination} for termination in GR or CGR is equivalent to $\chi_{t,I_t}(r_t')=1$, which we simply write as $\chi_t(r_t')=1$.

\subsubsection{Proof of Lemma~\ref{lem: general idea}}\label{subsection:proof_idea}

\textbf{Lemma~\ref{lem: general idea} (restated)} 
\textit{Consider resampling of $r_t'$ from $\mathcal{D}$ conditioned on $\nec_t$ until \eqref{cond_termination} is satisfied. Then, the number $M_t$ of resampling satisfies}
\begin{align}
\E_{r_t' \sim \mathcal{D}|\nec_t}[M_t|\hat{L}_t, I_t]=\frac{\sP[\nec_t|\hat{L}_t, I_t]}{w_{t,I_t}}.\n
\end{align}

\begin{proof2}
    Consider the arm-selection probability $w_{t,I_t}$ with the condition $\nec_t$. $w_{t,I_t}$ can be expressed as 
    \begin{align*}
        w_{t,I_t} &= \sP[\chi_t(r_t')=1|\hat{L}_t, I_t] \\
        &= \sP[\chi_t(r_t')=1|\nec_t,\hat{L}_t, I_t]\sP[\nec_t|\hat{L}_t, I_t]+\sP[\chi_t(r_t')=1|\nec_t^c,\hat{L}_t, I_t]\sP[\nec_t^c|\hat{L}_t, I_t]. \numberthis{\label{eq: conditioned on nec_t or not}}
    \end{align*}
    Note that $\nec_{t}$ is an arbitrary necessary condition for $\chi_t(r_t')=1$, which implies that
    \begin{equation*}
        \sP[\chi_t(r_t')=1|\nec_t^c,\hat{L}_t, I_t] = 0.
    \end{equation*}
    Therefore, from \eqref{eq: conditioned on nec_t or not} we have 
    \begin{equation}
        w_{t,I_t} = \sP[\chi_t(r_t')=1|\nec_t,\hat{L}_t, I_t]\sP[\nec_t|\hat{L}_t, I_t].
        \label{eq: conditioned on nec_t}
    \end{equation}
    Now we consider the expected number of resampling $M_t$. Recall that $r_t'$ is sampled from $\mathcal{D}$ conditioned on $\nec_t$ until \eqref{cond_termination} is satisfied, that is, $\chi_t(r_t')=1$. Then $M_t$ follows geometric distribution with probability mass function
    \begin{equation*}
        \sP[M_t=m|\hat{L}_t, I_t]=\qty(1-\sP[\chi_t(r_t')=1|\nec_t,\hat{L}_t, I_t])^{m-1}\sP[\chi_t(r_t')=1|\nec_t,\hat{L}_t, I_t].
    \end{equation*}
    Therefore, the expected number of resampling given $\hat{L}_t$ and $I_t$ is expressed as 
    \begin{align*}
        \E_{r_t' \sim \mathcal{D}|\nec_t}[M_t|\hat{L}_t, I_t] 
        &= \sP[\chi_t(r_t')=1|\nec_t,\hat{L}_t, I_t]\sum_{n=1}^{\infty}n\qty(1-\sP[\chi_t(r_t')=1|\nec_t,\hat{L}_t, I_t])^{n-1} \\
        &= \sP[\chi_t(r_t')=1|\nec_t,\hat{L}_t, I_t] /\qty(\sP[\chi_t(r_t')=1|\nec_t,\hat{L}_t, I_t])^2 \\
        &= 1/\sP[\chi_t(r_t')=1|\nec_t,\hat{L}_t, I_t]. \numberthis{\label{eq: expected number}}
    \end{align*}
    Combining \eqref{eq: conditioned on nec_t} and \eqref{eq: expected number}, we obtain
    \begin{align*}
        \E_{r_t' \sim \mathcal{D}|\nec_t}[M_t|\hat{L}_t, I_t]=\frac{\sP[\nec_t|\hat{L}_t, I_t]}{w_{t,I_t}}. \dqed
    \end{align*}
\end{proof2}

\subsubsection{Proof of Lemma~\ref{lem: complexity of CGR1}}\label{subsection:proof_cgr1}
\textbf{Lemma~\ref{lem: complexity of CGR1} (restated)} 
\textit{The sample $r_t'$ obtained by Algorithm~\ref{alg:CGR1} follows the conditional distribution of $\mathcal{D}$ given $\nec_t=\qty{r_{t,I_t}'=\max_{i:\sigma_i\leq\sigma_{I_t}}r_{t,i}'}$.}
\textit{In addition,}
\begin{align*}
\sP_{r_t' \sim \mathcal{D}}[\nec_t|\hat{L}_t,I_t]=1/\sigma_{I_t}
\end{align*}
\textit{and the number $M_t$ of resampling satisfies}
\begin{align*}
    \E_{r_t'\sim \mathcal{D}|\nec_t}[M_t|\hat{L}_t]\leq\log K+1.
\end{align*}

\begin{proof2}
    Let $\sP^*[\cdot]$ denote the probability distribution of $r_t'$ after the value-swapping operation. We have
    \begin{align*}
        &\sP^*\qty[\bigcap\nolimits_{i:\sigma_i\leq\sigma_{I_t}} \left\{r'_{t,i} \leq a_i\right\}\middle |\hat{L}_t, I_t] \\ 
        =& \sum_{j:\sigma_j\leq\sigma_{I_t}}\sP\qty[\bigcap\nolimits_{i:\sigma_i\leq\sigma_{I_t},i\notin\qty{j,I_t}} \left\{r'_{t,i} \leq a_i\right\},\, r'_{t,j} \leq a_{I_t},\, r'_{t,I_t} \leq a_j,\, r_{t,j}'=\max\limits_{i:\sigma_i\leq\sigma_{I_t}}r_{t,i}'\middle |\hat{L}_t, I_t] \\
        =& \sum_{j:\sigma_j\leq\sigma_{I_t}}\sP\qty[\bigcap\nolimits_{i: \sigma_i\leq\sigma_{I_t},i\notin\qty{j,I_t}} \left\{r'_{t,i} \leq a_i\right\},\, r'_{t,j} \leq a_{I_t},\, r'_{t,I_t} \leq a_j \middle | r_{t,j}'=\max\limits_{i:\sigma_i\leq\sigma_{I_t}}r_{t,i}',\hat{L}_t, I_t]\\
        &\hspace{8.6cm}
        \sP\qty[r_{t,j}'=\max\limits_{i:\sigma_i\leq\sigma_{I_t}}r_{t,i}'\middle |\hat{L}_t, I_t]. \numberthis{\label{eq: prob_swap}}
    \end{align*}
By the symmetry of $r_t'\in [0,\infty)^K$, we have
\begin{align}
\sP\qty[r_{t,j}'=\max\limits_{i:\sigma_i\leq\sigma_{I_t}}r_{t,i}'\middle |\hat{L}_t, I_t]
=\sP\qty[r_{t,I_t}'=\max\limits_{i:\sigma_i\leq\sigma_{I_t}}r_{t,i}'\middle |\hat{L}_t, I_t].
\label{max_prob}
\end{align}
for any $j$ suth that $\sigma_j\leq\sigma_{I_t}$.
Then we have
\begin{align}
1
&=
\sP\qty[\bigcup_{j: \sigma_j\leq\sigma_{I_t}}
\{r_{t,j}'=\max\limits_{i:\sigma_i\leq\sigma_{I_t}}r_{t,i}'\}\middle |\hat{L}_t, I_t]
\nn
&=
\sum_{j: \sigma_j\leq\sigma_{I_t}}
\sP\qty[
r_{t,j}'=\max\limits_{i:\sigma_i\leq\sigma_{I_t}}r_{t,i}'\middle |\hat{L}_t, I_t]
\nn
&=
\sigma_{I_t}
\sP\qty[
r_{t,I_t}'=\max\limits_{i:\sigma_i\leq\sigma_{I_t}}r_{t,i}'\middle |\hat{L}_t, I_t],\n
\end{align}
which means that \eqref{max_prob} is equal to $1/\sigma_{I_t}$.
%
%    For any $j$ that satisfies $\sigma_j\leq\sigma_{I_t}$, we have 
%    \begin{align*}
%        \sP\qty[r_{t,j}'=\max_{\sigma_i\leq\sigma_{I_t}}r_{t,i}'\middle |\hat{L}_t, I_t] 
%        &= \int_0^{\infty} f(z)F^{\sigma_{I_t}-1}(z) \,dz \\
%        &= \int_0^{\infty} F^{\sigma_{I_t}-1}(z) \,dF(z) \\
%        &= 1/\sigma_{I_t}.
%        \numberthis{\label{eq: CGRI nec_prob}}
%    \end{align*}
    Therefore, from \eqref{eq: prob_swap} we have 
    \begin{multline}
        \sP^*\qty[\bigcap\nolimits_{i: \sigma_i\leq\sigma_{I_t}} \left\{r'_{t,i} \leq a_i\right\}\middle |\hat{L}_t, I_t] = \\
        \frac{1}{\sigma_{I_t}}\sum_{j:\sigma_j\leq\sigma_{I_t}}\sP\qty[\bigcap\nolimits_{i: \sigma_i\leq\sigma_{I_t},i\notin\qty{j,I_t}} \left\{r'_{t,i} \leq a_i\right\},\, r'_{t,j} \leq a_{I_t},\, r'_{t,I_t} \leq a_j \middle | r_{t,j}'=\max\limits_{i:\sigma_i\leq\sigma_{I_t}}r_{t,i}',\hat{L}_t, I_t].
        \label{eq: prob_swap2}
    \end{multline}
    By symmetry, each probability term on RHS of \eqref{eq: prob_swap2} is equal. Therefore, we have
    \begin{equation*}
        \sP^*\qty[\bigcap\nolimits_{i: \sigma_i\leq\sigma_{I_t}} \left\{r'_{t,i} \leq a_i\right\}\middle |\hat{L}_t, I_t] = \sP\qty[\bigcap\nolimits_{i: \sigma_i\leq\sigma_{I_t}} \left\{r'_{t,i} \leq a_i\right\} \middle | \nec_t, \hat{L}_t, I_t],
    \end{equation*}
    which means that CGR~I samples $r_t'$ from the conditional distribution of $\mathcal{D}$ given $\nec_t=\qty{r_{t,I_t}'=\max_{i:\sigma_i\leq\sigma_{I_t}}r_{t,i}'}$.

    Combining this fact with Lemma~\ref{lem: general idea}, for CGR I we have
    \begin{equation*}
        \E_{r_t' \sim \mathcal{D}|\nec_t}[M_t|\hat{L}_t, I_t]=\frac{1}{\sigma_{I_t}w_{t,I_t}}.
    \end{equation*}
    Then, the expected number of resampling given $\hat{L}_t$ in CGR~I is bounded by 
    \begin{align*}
        \E_{r_t' \sim \mathcal{D}|\nec_t}[M_t|\hat{L}_t] &= \sum_{i=1}^K \sP[I_t=i|\hat{L}_t]\E_{r_t' \sim \mathcal{D}|\nec_t}[M_t|\hat{L}_t, I_t=i] \\
        &= \sum_{i=1}^K w_{t,i}\cdot\frac{1}{\sigma_i w_{t,i}} \\
        &\leq 1+\int_{1}^{K} \frac{1}{x} \,dx  \\
        &= \log K + 1. \dqed
    \end{align*}
\end{proof2}

\subsubsection{Proof of Lemma~\ref{lem: complexity of CGR2}}\label{subsection:proof_cgr2}
\textbf{Lemma~\ref{lem: complexity of CGR2} (restated)} 
\textit{In CGR~II, $r_t'$ follows the conditional distribution of $\mathcal{D}$ given $\nec_t$ in \eqref{A_CGR2}. In addition,}
\begin{align*}
    \sP_{r_t' \sim \mathcal{D}}[\nec_t|\hat{L}_t,I_t]=\qty(1-F^{\sigma_{I_t}}(\eta_t\underline{\hat{L}}_{t,I_t}))/\sigma_{I_t}
\end{align*} 
\textit{and the number $M_t$ of resampling satisfies}
\begin{align}
    \E_{r_t' \sim \mathcal{D}}[M_t|\hat{L}_t]\leq\log K+1.\n
\end{align}
\textit{Besides, if $\mathcal{D}$ is the Fr\'{e}chet distribution with shape $2$, then $M_t$ satisfies}
\begin{align*}
    \E_{r_t' \sim \mathcal{D}|\nec_t}[M_t|\hat{L}_t,I_t]\leq K\lor 4.
\end{align*}

\begin{proof}
    Recall that $\mathcal{D}_{\sigma_{I_t}}$ be the distribution of the maximum of $\sigma_{I_t}$ i.i.d.~samples from $\mathcal{D}$, whose cumulative distribution function is given by 
    \begin{align*}
        F^{\sigma_{I_t}}(x)=(F(x))^{\sigma_{I_t}}.
    \end{align*}
    Then, the probability density function of $\mathcal{D}_{\sigma_{I_t}}$ is expressed as 
    \begin{align*}
        f^{\sigma_{I_t}}(x)=\begin{cases}
            \sigma_{I_t}f(x)F^{\sigma_{I_t}-1}(x), & x\in\left[0, \infty\right), \\
            0, & \text{otherwise}.
        \end{cases}
    \end{align*}
    Since $r_{t,I_t}'$ is sampled from the truncated distribution of $\mathcal{D}_{\sigma_{I_t}}$ with support $\left[\eta_t\underline{\hat{L}}_{t,I_t}, \infty\right)$, it follows the probability density function 
    \begin{equation*}
        f_{I_t}(x; \eta_t\underline{\hat{L}}_{t,I_t})=
        \begin{cases}
            \sigma_{I_t}f(x)F^{\sigma_{I_t}-1}(x)/\qty(1-F^{\sigma_{I_t}}\qty(\eta_t\underline{\hat{L}}_{t,I_t})), & x\in\left[\eta_t\underline{\hat{L}}_{t,I_t}, \infty\right), \\
            0, & \text{otherwise}.
        \end{cases}
    \end{equation*}
    For any $i\neq I_t$ that satisfies $\sigma_i\leq\sigma_{I_t}$, given $r_{t,I_t}'=a_{I_t}$, $r_{t,i}'$ is sampled i.i.d. from the truncated distribution of $\mathcal{D}$ with support $\left[0, a_{I_t}\right]$. Therefore, it follows the probability density function 
    \begin{equation*}
        f_i(x; a_{I_t})=
        \begin{cases}
            f(x)/F(a_{I_t}), & x\in\left[0, a_{I_t}\right], \\
            0, & \text{otherwise}.
        \end{cases}
    \end{equation*}
    Given $\underline{\hat{L}}_{t,I_t}$, the joint probability density of $r_{t,i}'$ for $i$ satisfying $\sigma_i\leq\sigma_{I_t}$ is denoted by $f^K(\cdot; \eta_t\underline{\hat{L}}_{t,I_t})$. 
    Then, for $a=\qty(a_1,a_2,\dots, a_K)$ we have
    \begin{multline*}
        f^K(a; \eta_t\underline{\hat{L}}_{t,I_t})
        = f_{I_t}(a_{I_t}; \eta_t\underline{\hat{L}}_{t,I_t})\prod_{i:\sigma_i\leq\sigma_{I_t},i\neq I_t}f_i(a_i; a_{I_t}) \\
        =
        \begin{cases}
            \sigma_{I_t}\prod_{i:\sigma_i\leq\sigma_{I_t}}f(a_i)/\qty(1-F^{\sigma_{I_t}}\qty(\eta_t\underline{\hat{L}}_{t,I_t})), & \text{if } a_{I_t}=\max_{i:\sigma_i\leq\sigma_{I_t}}a_i, a_{I_t}\geq\eta_t\underline{\hat{L}}_{t,I_t}, \\
            0, & \text{otherwise}.
        \end{cases}
    \end{multline*}
    Let $\sP^*[\cdot]$ denote the probability distribution of $r_t'$ under the resampling method in CGR~II. We have
    \begin{align*}
        \lefteqn{\sP^*\qty[\bigcap\nolimits_{i:\sigma_i\leq\sigma_{I_t}} \left\{r'_{t,i} \leq a_i\right\}\middle |\hat{L}_t, I_t]}\\
        &= \underbrace{\int_0^{a_i}\cdots\int_0^{a_{i'}}}_{\sigma_{I_t} \text{integrals}}f_{I_t}(z_{I_t}; \eta_t\underline{\hat{L}}_{t,I_t}) \,dz_{I_t} \prod_{i: \sigma_i\leq\sigma_{I_t},i\neq I_t}f_i(z_i; z_{I_t}) \,dz_i \numberthis{\label{eq: distinguish}}\\
        &= \frac{\sigma_{I_t}}{1-F^{\sigma_{I_t}}\qty(\eta_t\underline{\hat{L}}_{t,I_t})}\int_{\eta_t\underline{\hat{L}}_{t,I_t}}^{a_{I_t}}f(z)\prod_{i: \sigma_i\leq\sigma_{I_t},i\neq I_t} F(a_i\wedge z) \,dz \\
        &= \frac{\sigma_{I_t}}{1-F^{\sigma_{I_t}}\qty(\eta_t\underline{\hat{L}}_{t,I_t})}\sP\qty[\bigcap\nolimits_{i:\sigma_i\leq\sigma_{I_t}} \left\{r'_{t,i} \leq a_i\right\}, \nec_t\middle |\hat{L}_t, I_t].
        \numberthis{\label{eq: CGRII sample_prob}}
    \end{align*}
    Here, $i'$ in \eqref{eq: distinguish} is used as a distinct index from $i$, satisfying the constraint $\sigma_{i'}\leq\sigma_{I_t}$. Now we consider the probability of $\nec_t=\qty{r_{t,I_t}'=\max_{i:\sigma_i\leq\sigma_{I_t}}r_{t,i}',r_{t,I_t}'\geq\eta_t\underline{\hat{L}}_{t,I_t}}$ given $\hat{L}_t$ and $I_t$, which is expressed as 
    \begin{align*}
        \sP[\nec_t|\hat{L}_t, I_t]
        &=\sP\qty[r_{t,I_t}'=\max_{i:\sigma_i\leq\sigma_{I_t}}r_{t,i}',r_{t,I_t}'\geq\eta_t\underline{\hat{L}}_{t,I_t}\middle |\hat{L}_t, I_t] \\
        &= \int_{\eta_t\underline{\hat{L}}_{t,I_t}}^{\infty} f(z)F^{\sigma_{I_t}-1}(z) \,dz \\
        &= \int_{\eta_t\underline{\hat{L}}_{t,I_t}}^{\infty} F^{\sigma_{I_t}-1}(z) \,dF(z) \\
        &= \qty(1-F^{\sigma_{I_t}}\qty(\eta_t\underline{\hat{L}}_{t,I_t}))/\sigma_{I_t}.
        \numberthis{\label{eq: CGRII nec_prob}}
    \end{align*}
    Combining \eqref{eq: CGRII sample_prob} and \eqref{eq: CGRII nec_prob}, we have
    \begin{align*}
        \sP^*\qty[\bigcap\nolimits_{i:\sigma_i\leq\sigma_{I_t}} \left\{r'_{t,i} \leq a_i\right\}\middle |\hat{L}_t, I_t] 
        &= \frac{\sP\qty[\bigcap\nolimits_{i:\sigma_i\leq\sigma_{I_t}} \left\{r'_{t,i} \leq a_i\right\}, \nec_t\middle |\hat{L}_t, I_t]}{\sP[\nec_t|\hat{L}_t, I_t]} \\
        &= \sP\qty[\bigcap\nolimits_{i:\sigma_i\leq\sigma_{I_t}} \left\{r'_{t,i} \leq a_i\right\}\middle |\nec_t, \hat{L}_t, I_t],
    \end{align*}
    which means that CGR~II samples $r_t'$ from $\mathcal{D}$ conditioned on
    \begin{equation*}
        \nec_t=\qty{r_{t,I_t}'=\max_{i:\sigma_i\leq\sigma_{I_t}}r_{t,i}',r_{t,I_t}'\geq\eta_t\underline{\hat{L}}_{t,I_t}}.
    \end{equation*}

    The value of $\sP[\nec_t|\hat{L}_t, I_t]$ is given in \eqref{eq: CGRII nec_prob}. According to Lemma~\ref{lem: general idea}, for CGR~II, $M_t$ satisfies
    \begin{equation*}
        \E_{r_t' \sim \mathcal{D}|\nec_t}[M_t|\hat{L}_t, I_t]\leq\frac{1-F^{\sigma_{I_t}}\qty(\eta_t\underline{\hat{L}}_{t,I_t})}{\sigma_{I_t}w_{t,I_t}},
    \end{equation*}
    where the equality holds if and only if $G_t=\infty$.
    
    Then, the expected number of resampling given $\hat{L}_t$ in CGR~II is bounded by 
    \begin{align*}
        \E_{r_t' \sim \mathcal{D}|\nec_t}[M_t|\hat{L}_t] &= \sum_{i=1}^K \sP[I_t=i|\hat{L}_t]\E_{r_t' \sim \mathcal{D}|\nec_t}[M_t|\hat{L}_t, I_t=i] \\
        &\leq \sum_{i=1}^K w_{t,i}\cdot\frac{1-F^{\sigma_{I_t}}\qty(\eta_t\underline{\hat{L}}_{t,I_t})}{\sigma_i w_{t,i}} \\
        &\leq 1+\int_{1}^{K} \frac{1}{x} \,dx  \\
        &= \log K + 1.
    \end{align*}
    According to Lemma~\ref{lem: general idea}, we have
    $\E_{r_t' \sim \mathcal{D}|\nec_t}[M_t|\hat{L}_t, I_t]=\sP[\nec_t|\hat{L}_t, I_t]/w_{t,I_t}$.
    Therefore, to prove $\E_{r_t' \sim \mathcal{D}|\nec_t}[M_t|\hat{L}_t, I_t=i]\allowbreak\leq K\lor 4$, we only need to show $\sP[\gA_{t}|\hat{L}_t, I_t=i]/w_{t,i}\leq K\lor 4$. See the proof in Section~\ref{subsection: B2}.
\end{proof}

\subsection{Proofs of regret analysis}\label{append_regret}
In this section, we provide the proofs for Lemmas~\ref{lem: regret decomposition} and~\ref{lem: additive term by CGR biased}
on the bias and the regret of CGR II-biased.

\subsubsection{Proof of Lemma~\ref{lem: regret decomposition}}
Before presenting the proof of Lemma~\ref{lem: regret decomposition}, we first provide an explicit expression for the expectation of the loss estimates generated by CGR II, similar to the approach used in Lemma 4 of \citet{JMLR:v17:15-091}.
\begin{lemma}\label{lem: bias of loss}
    For all $i\in [K]$ and $t$ such that $w_{t,i} > 0$, the loss estimates of CGR II satisfies that
    \begin{equation*}
        \E\qty[\hat{\ell}_{t,i} \middle| \hat{L}_t] = \qty(1-\qty(1-\frac{w_{t,i}}{\sP(\gA_{t}|\hat{L}_t, I_t=i)})^{G_{t}}) \ell_{t,i}.
    \end{equation*}
\end{lemma}
\begin{proof}
    In CGR II, the probability of pulling an arm $i$, denoted by $q_{t,i}$, is given as 
% \honda{Please unify whether we ``play'' arms or ``pull'' arms.}
    \begin{multline*}
       q_{t,i}= \sP\qty(\chi_{t,i}(r_{t,i}') = 1 \middle | \hat{L}_t) = \sP\qty(\chi_{t,i}(r_{t,i}') = 1 \middle | \hat{L}_t, I_t=i, \gA_{t}) \sP(\gA_{t}| \hat{L}_t, I_t=i)\\
        =  \frac{w_{t,i}}{\sP(\gA_{t}| \hat{L}_t, I_t=i)},
    \end{multline*}
    where the last equality follows from (\ref{eq: conditioned on nec_t}).
    Let $M_{t,i}$ denote the number of resampling for an arm $i$ at round $t$.
    % Denote $q= q_{t,i}$ and $G=G_{t,i}$ for brevity.
    Then, we have
    \begin{align*}
        \E_{r_t' \sim \mathcal{D}|\nec_t}[M_{t,i} | \hat{L}_t] &= \sum_{k=1}^\infty k(1-q_{t,i})^{k-1} q_{t,i} - \sum_{k=G_{t}}^\infty (k-G_t)(1-q_{t,i})^{k-1}q_{t,i} \\
        &= \sum_{k=1}^\infty k(1-q_{t,i})^{k-1}q_{t,i} - (1-q_{t,i})^{G_t} \sum_{k=G_t}^\infty (k-G_t)(1-q_{t,i})^{k-G_t-1} q_{t,i} \\
        &= (1-(1-q_{t,i})^{G_t}) \sum_{k=1}^\infty k(1-q_{t,i})^{k-1}q_{t,i} = \frac{1-(1-q_{t,i})^{G_t}}{q_{t,i}}.
    \end{align*}
    By definition of $\hat{\ell}_{t,i}$, we have
    \begin{equation*}
        \E\qty[ \hat{\ell}_{t,i} \middle | \hat{L}_t] = q_{t,i} \ell_{t,i} \E\qty[M_{t,i} \middle | \hat{L}_t] = (1-(1-q_{t,i})^{G_{t,i}}) \ell_{t,i}. \qedhere
    \end{equation*}
% \honda{I wonder that currently the benefit to use $q, G$ is limited.
% How about also rewriting the final equation using them (or just writing using $q_{t,i},G_{t,i}$ everywhere)?}
\end{proof}


\noindent\textbf{Lemma~\ref{lem: regret decomposition} (restated)} 
\textit{The expected regret of FTPL with CGR II satisfies}
\begin{equation*}
    \gR(T) \leq \sum_{t=1}^T \mathbb{E}\qty[\inp{\hat{\ell}_t}{w_t - e_{i^*}}]
    + \sum_{t=1}^T \sum_{i\in [K]} \E\qty[w_{t,i} \qty(1-\frac{w_{t,i}}{\sP(\gA_{t}|\hat{L}_t, I_t=i)})^{G_t}].
\end{equation*}
\begin{proof}
    The overall proof is a simple reduction of the results in combinatorial semi-bandits~\citep[Lemma~5]{JMLR:v17:15-091} to MAB.
     By definition, we have for any $i\in [K]$
    \begin{align*}
       \E\qty[ w_{t,i}\hat{\ell}_{t,i} \middle| \hat{L}_t] &=  w_{t,i} \E\qty[\hat{\ell}_{t,i} \middle| \hat{L}_t] \\
        &=  w_{t,i} \qty(1-\qty(1-\frac{w_{t,i}}{\sP(\gA_{t}|\hat{L}_t, I_t=i)})^{G_{t,i}})    \ell_{t,i} \tag{by Lemma~\ref{lem: bias of loss}} \\
        &=  w_{t,i} \ell_{t,i} -   w_{t,i}\qty(1-\frac{w_{t,i}}{\sP(\gA_{t}|\hat{L}_t, I_t=i)})^{G_{t,i}}\ell_{t,i} \\
        & \geq w_{t,i} \ell_{t,i} -   w_{t,i}\qty(1-\frac{w_{t,i}}{\sP(\gA_{t}|\hat{L}_t, I_t=i)})^{G_{t,i}}.
        \numberthis{\label{eq: reduction 1}}
    \end{align*}
    Then, by definition of $\gR(T)$, we have
\begin{align*}
    \gR(T) &= \sum_{t=1}^T \E\qty[\inp{\ell_t}{w_t - e_{i^*}}] \\
    &\leq \sum_{t=1}^T \E\qty[\inp{\hat{\ell}_t}{w_t} - \inp{\ell_t}{e_{i^*}}] +\sum_{t=1}^T \sum_{i \in [K]}  \E\qty[w_{t,i}\qty(1-\frac{w_{t,i}}{\sP(\gA_{t}|\hat{L}_t, I_t=i)})^{G_{t,i}}]\tag{by (\ref{eq: reduction 1})},
\end{align*}
which concludes the proof.
\end{proof}

\subsubsection{Proof of Lemma~\ref{lem: additive term by CGR biased}}\label{subsection: B2}
\textbf{Lemma~\ref{lem: additive term by CGR biased} (restated)} 
\textit{When $\mathcal{D}$ is the Fr\'{e}chet distribution with shape $2$ and $G_t=(K\lor 4)\log t$, FTPL with CGR II-biased satisfies}
\begin{equation*}
    \sum_{t=1}^T \sum_{i\in [K]}\E\qty[w_{t,i} \qty(1-\frac{w_{t,i}}{\sP(\gA_{t}|\hat{L}_t, I_t=i)})^{G_t}] \leq \log T.
\end{equation*}
\begin{proof}
    For notational simplicity, we denote $\sP(\gA_{t}|\hat{L}_t, I_t=i)$ by $\sP(\gA_{t,i})$ in this proof.
    Since $(1-x)^a \leq e^{-ax}$, it suffices to show that for any $t \in \sN$ and $i \in [K]$
    \begin{equation*}
       \exp(-\frac{w_{t,i}}{\sP(\gA_{t,i})}G_t) \leq \frac{1}{t}.
    \end{equation*}
% \honda{I rephrased a bit since it is nontrivial whether it is really ``necessary''.}
To this end, we
%    Therefore, it is necessary to
address the term $\frac{\sP(\gA_{t,i})}{w_{t,i}}$ and set an appropriate $G_t$ to establish a tight regret upper bound.

    By definition of $\gA_{t,i}$, when $\sigma_i$ denotes the number of arms such that $\hat{L}_{t,j} \leq \hat{L}_{t,i}$ it holds that
    \begin{equation*}
        \sP(\gA_{t,i}) = \int_{\ul_{t,i}}^\infty f(z) F^{\sigma_i - 1}(z) \dd z.
    \end{equation*}
    When $\gD$ is
%given as
Fr\'{e}chet distribution with shape $2$, whose density and distribution functions are given as
    \begin{equation}\label{eq: Frechet density}
        f(x) = \frac{2 e^{-\frac{1}{x^2}}}{x^3}, \qq{and} F(x) = e^{-\frac{1}{x^2}},
    \end{equation}
    it holds that
    \begin{align*}
        \sP(\gA_{t,i}) &= \int_{\ul_{t,i}}^\infty \frac{2}{z^{3}}\exp(-\frac{\sigma_i}{z^2} ) \dd z \\
        &= \int_0^{\frac{\sigma_i}{\ul_{t,i}^2}} \frac{1}{\sigma_i} e^{-t} \dd t \tag{by $t=\sigma_i/z^2$} \\
        &= \frac{1-\exp(-\frac{\sigma_i}{\ul_{t,i}^2})}{\sigma_i}. \numberthis{\label{eq: expression of prob. A}}
    \end{align*}
Now,
%Then,
let us consider the lower bound on $w_{t,i}$.
    When $z>0$ and $\ul_{t,j} \geq \ul_{t,i}$, i.e., $\sigma_j \geq \sigma_i$ we have
    \begin{equation*}
        \frac{1}{(z+\ul_{t,j})^2} \leq \frac{1}{(z+\ul_{t,i})^2}.
    \end{equation*}
    On the other hand, when $z \geq \ul_{t,i}$ and $\ul_{t,j} \leq \ul_{t,i}$, we have
    \begin{equation*}
        \frac{1}{(z+\ul_{t,j})^2} < \frac{1}{\qty(\frac{z+\ul_{t,i}}{2}+\ul_{t,j})^2} < \frac{4}{(z+\ul_{t,i})^2}.
    \end{equation*}
    Therefore, we obtain
    \begin{align*}
        w_{t,i} &= \int_0^\infty \frac{2}{(z+\ul_{t,i})^3} \exp(-\sum_{j \in [K]} \frac{1}{(z+\ul_{t,j})^2}) \dd z \\
        &\geq \int_{\ul_{t,i}}^\infty \frac{2}{(z+\ul_{t,i})^3} \exp(-\sum_{j \in [K]} \frac{1}{(z+\ul_{t,j})^2}) \dd z \\
        &\geq \int_{\ul_{t,i}}^\infty \frac{2}{(z+\ul_{t,i})^3} \exp(- \frac{4(\sigma_i-1) + K-\sigma_i +1}{(z+\ul_{t,i})^2}) \dd z  \\
        &= \frac{1-\exp(-\frac{3(\sigma_i -1)+K}{4\ul_{t,i}^2})}{3(\sigma_i-1)+K}. \numberthis{\label{eq: lower bound on w}}
    \end{align*}
    Let $h(x)=(1-e^{-x})/x$.
    Then, by combining (\ref{eq: expression of prob. A}) and (\ref{eq: lower bound on w}), we obtain
        \begin{align*}
       \frac{\sP(\gA_{t,i})}{w_{t,i}}
        &\leq \frac{1-\exp(-\frac{\sigma_i}{\ul_{t,i}^2})}{\sigma_i} \frac{ 3(\sigma_i-1)+K}{1-\exp(-\frac{3(\sigma_i -1)+K}{4\ul_{t,i}^2})} \\
        &\leq \frac{4h(x)}{h(y)}, \numberthis{\label{eq: P/w ratio}}
    \end{align*}
    where $x=\sigma_i/\ul_{t,i}^2$ and $y=\frac{3(\sigma_i-1)+K}{4\ul_{t,i}^2}$.
    Except the case $\sigma_i=K=2$, we can see 
    \begin{equation*}
        \frac{y}{x} = \frac{3(\sigma_i-1)+K}{4\sigma_i} \leq \frac{K}{4}
    \end{equation*}
    and therefore
    \begin{align*}
    \frac{\sP(\gA_{t,i})}{w_{t,i}} &\le \sup_{x,y>0: y\le Kx/4} \frac{4h(x)}{h(y)}
    \\
    &= \sup_{x>0} \frac{4h(x)}{h(Kx/4)},
    \end{align*}
    since $h(y)$ is decreasing.
    For $K\leq 4$, it is bounded from above by $4$ since $h(x) \leq h(Kx/4)$ by $x\geq Kx/4$.
    For $K \geq 4$, we have
    \begin{align*}
        \sup_{x>0}
        \frac{4h(x)}{h(Kx/4)}
        &=
        K\sup_{x>0}
        \frac{1-e^{-x}}{1-e^{-Kx/4}} \numberthis{\label{eq: h ratio}}\\
        &=
        K\sup_{x>0}
        \left(1-
        \frac{e^{-x}-e^{-Kx/4}}{1-e^{-Kx/4}}
        \right)
        \le K.
    \end{align*}
    Finally, when $\sigma_i = K=2$, $y/x=5/8$ holds.
    Since $h(y)$ is decreasing, we again obtain
    \begin{equation*}
         \frac{\sP(\gA_{t,i})}{w_{t,i}} \le \sup_{x>0} \frac{4h(x)}{h(5x/8)} \leq 4.
    \end{equation*}
    Hence, for all $t \in \sN$ and $i \in [K]$, we obtain
    \begin{align*}
        \frac{\sP(\gA_{t,i})}{w_{t,i}}
        &\le
        K\lor 4.
    \end{align*}
    % Therefore, for all $t \in \sN$ and $i \in [K]$, we obtain
    % \begin{equation*}
    %     \frac{\sP(\gA_{t,i})}{w_{t,i}} \leq K+7 \implies \frac{w_{t,i}}{\sP(\gA_{t,i})} \geq \frac{1}{K+7}.
    % \end{equation*}
    This implies that for any $G_t >0$
    \begin{align*}
        \sum_i w_{t,i}\exp(-\frac{w_{t,i}}{\sP(\gA_{t,i})}G_t)  &\leq \sum_i w_{t,i}\exp(-\frac{G_t}{K \lor 4}) \\
        &= \exp(-\frac{G_t}{K \lor 4}). \tag{by $ \sum_i w_{t,i} = 1$}
    \end{align*}
    Therefore, setting $G_t = (K\lor 4) \log t$ concludes the proof.
\end{proof} 



        
%     By combining (\ref{eq: expression of prob. A}) and (\ref{eq: lower bound on w}), we obtain
%     \begin{align*}
%         \frac{\sP(\gA_{t,i})}{w_{t,i}}
% &\leq \frac{1-\exp(-\frac{\sigma_i}{\ul_{t,i}^2})}{\sigma_i} \frac{ 3(\sigma_i-1)+K}{1-\exp(-\frac{3(\sigma_i -1)+K}{4\ul_{t,i}^2})} \\
%         &\leq \frac{1-\exp(-\frac{\sigma_i}{\ul_{t,i}^2})}{\sigma_i} \frac{4\ul_{t,i}^2 + 3(\sigma_i -1)+K}{3(\sigma_i-1)+K}  (3(\sigma_i-1)+K) \numberthis{\label{eq: middle eq for the ratio 1}} \\
%         & = \frac{4\ul_{t,i}^2 + 3(\sigma_i -1)+K}{\sigma_i} \qty(1-\exp(-\frac{\sigma_i}{\ul_{t,i}^2})),
%     \end{align*}
%     where (\ref{eq: middle eq for the ratio 1}) follows from $\frac{x}{1+x} < 1-e^{-x}$ for $x>-1$.

% \honda{I feel that this kind of discussion is tighter:
% Let $f(x)=(1-e^{-x})/x$.
% Then,
% \begin{align}
% \frac{\sP(\gA_{t,i})}{w_{t,i}}
% &\leq \frac{1-\exp(-\frac{\sigma_i}{\ul_{t,i}^2})}{\sigma_i} \frac{ 3(\sigma_i-1)+K}{1-\exp(-\frac{3(\sigma_i -1)+K}{4\ul_{t,i}^2})}
% \nn
% &=
% \frac{4f(x)}{f(y)}
% \n
% \end{align}
% for $x=\sigma_i/\ul_{t,i}^2$ and $y=\frac{3(\sigma_i-1)+K}{4\ul_{t,i}^2}$.
% Except the case $\sigma_i=K=2$ we can see $y/x\le K/4$
% and therefore
% \begin{align}
% \frac{\sP(\gA_{t,i})}{w_{t,i}}
% &\le
% \sup_{x,y>0: y\le Kx/4}
% \frac{4f(x)}{f(y)}
% \nn
% &=
% \sup_{x>0}
% \frac{4f(x)}{f(Kx/4)}
% \n
% \end{align}
% since $f(y)$ is decreasing.
% For $K\le 4$, it is bounded from above by 4
% since $f(x)\le f(Kx/4)$ by $x\ge Kx/4$.
% For $K\ge 4$ we have
% \begin{align}
% \sup_{x>0}
% \frac{4f(x)}{f(Kx/4)}
% &=
% K\sup_{x>0}
% \frac{1-e^{-x}}{1-e^{-Kx/4}}
% \nn
% &=
% K\sup_{x>0}
% \left(1-
% \frac{e^{-x}-e^{-Kx/4}}{1-e^{-Kx/4}}
% \right)
% \le K.
% \n
% \end{align}
% Combining them we have
% \begin{align}
% \frac{\sP(\gA_{t,i})}{w_{t,i}}
% &\le
% K\lor 4.
% \n
% \end{align}
% We can confirm that this becomes true also for $\sigma_i=K=2$ where $y/x=5/8$.
%\begin{itemize}
%\item To use
%\begin{align}
%\lefteqn{
%\frac{4\ul_{t,i}^2 + 3(\sigma_i -1)+K}{\sigma_i} \qty(1-\exp(-\frac{\sigma_i}{\ul_{t,i}^2})),
%}\\
%&\le
%\left(
%\frac{4\ul_{t,i}^2}{\sigma_i}+K
%\right)
%\qty(1-\exp(-\frac{\sigma_i}{\ul_{t,i}^2})),
%\end{align}
%(it is unclear for $K=2$ but we can confirm it manually for $\sigma_i=1,2$.)
%\item To change the case analysis from $\ul_{t,i}^2 \leq \sigma_i$ to $c\ul_{t,i}^2 \leq \sigma_i$ for some constant $c\ge 1$.
%\item To apply relation $1-e^{-x}\le 1-e^{-c}+e^{-c}(x-c)$ for the case analysis of $c\ul_{t,i}^2 \le \sigma_i$.
%Then the maximum would be when $\ul_{t,i}^2 = c\sigma_i$ under reasonable $c$.
%The optimization of the parameters could be done by a computer.
%\end{itemize}
% }

% When $\ul_{t,i}^2 \leq \sigma_i$, we have
% \begin{align*}
%      \frac{4\ul_{t,i}^2 + 3(\sigma_i -1)+K}{\sigma_i} \qty(1-\exp(-\frac{\sigma_i}{\ul_{t,i}^2})) &\leq \frac{4\ul_{t,i}^2 + 3(\sigma_i -1)+K}{\sigma_i} \\
%      &\leq \frac{7\sigma_i + K }{\sigma_i} \leq K+7.
% \end{align*}
% On the other hand, when $\ul_{t,i}^2 > \sigma_i$, we have
% \begin{align*}
%     \frac{4\ul_{t,i}^2 + 3(\sigma_i -1)+K}{\sigma_i} \qty(1-\exp(-\frac{\sigma_i}{\ul_{t,i}^2})) &\leq \frac{4\ul_{t,i}^2 + 3(\sigma_i -1)+K}{\sigma_i} \frac{\sigma_i}{\ul_{t,i}^2} \numberthis{\label{eq: middle eq for the ratio 2}} \\
%     &\leq 4 + \frac{3(\sigma_i-1)+K}{\ul_{t,i}^2} \\
%     &\leq 4+ \frac{3(\sigma_i-1)+K}{\sigma_i} \leq K+7,
% \end{align*}
% where (\ref{eq: middle eq for the ratio 2}) follows from $1-e^{-x} < x$ for $x>0$.

\subsubsection{Proof of Theorem~\ref{thm_biased}}
\textbf{Theorem~\ref{thm_biased} (restated)}
\textit{FTPL with CGR II-biased with learning rate $\eta_t = c/\sqrt{t}$ for $c>0$ and maximum number of resampling $G_{t}=(K\lor 4)\log t$ satisfies that
    \begin{equation*}
        \mathcal{R}(T) \leq \begin{cases}
            O\qty (\sqrt{KT})  &\text{in adversarial bandits,} \\
            O\qty(\sum_{i \ne i^*}\frac{\log T}{\Delta_i})& \text{in stochastic bandits,}
        \end{cases}
    \end{equation*}
if the perturbation distribution $\gD$ is the Fr\'{e}chet distribution with shape $2$.}
\begin{proof}
    The regret of FTPL with CGR II-unbiased with $G_t =(K\lor 4)\log t $ satisfies
    \begin{align*}
        \gR(T) &\leq \sum_{t=1}^T \mathbb{E}\qty[\inp{\hat{\ell}_t}{w_t - e_{i^*}}]
    + \sum_{t=1}^T \sum_{i\in [K]} \E\qty[w_{t,i} \qty(1-\frac{w_{t,i}}{\sP(\gA_{t}|\hat{L}_t, I_t=i)})^{G_t}] \tag{by Lemma~\ref{lem: regret decomposition}} \\
    &\leq \sum_{t=1}^T \mathbb{E}\qty[\inp{\hat{\ell}_t}{w_t - e_{i^*}}]
    + \log T .\tag{by Lemma~\ref{lem: additive term by CGR biased}}
    \end{align*}
    Let $\hat{\ell}_t^{\mathrm{CGR}}$ and $\hat{\ell}_t^{\mathrm{GR}}$ denote the IW estimators of CGR II-unbiased and GR, respectively.
    By construction, $\E[\hat{\ell}_t^{\mathrm{CGR}}] \leq \E[\hat{\ell}_t^{\mathrm{GR}}]$ holds, allowing the first term on the RHS to be directly bounded by the results for GR in \citet{pmlr-v201-honda23a}.
    Further details on adversarial regret can be found in Lemma~\ref{lem: main rslt of hnd} in Section~\ref{app: CGR II}.
    For the stochastic setting, as the analysis remains identical to the previous one, we refer readers to \citet[Theorem 2]{pmlr-v201-honda23a} for further details.
\end{proof}

\subsection{Improved regret analysis of CGR II}\label{app: CGR II}
In this section, we formalize Remark~\ref{thm: informal} by presenting an improved analysis of the adversarial regret for FTPL with CGR II when the perturbation $\gD$ follows the Fréchet distribution with shape parameter 2, as detailed in Theorem~\ref{thm: formal}.
To understand how the transition from GR to CGR II affects the regret analysis, we first introduce the necessary notations and summarize relevant prior results.
With a slight abuse of notation, we define a function $\phi_i$ for $i\in [K]$ by
\begin{equation*}
    \phi_i(\lambda) = \int_0^\infty \frac{2}{(z+\ul_i)^3} \exp(-\sum_{j=1}^{K} \frac{1}{(z+\ul_i)^2}) \dd z.
\end{equation*}
When $\gD$ is Fr\'echet distribution with shape $2$, whose density and distribution functions are given in (\ref{eq: Frechet density}), one can see that $\phi_i(\eta_t \hat{\uL}_{t}) = w_{t,i}$ holds.
Let $\phi = (\phi_1, \ldots, \phi_K)$ denote arm-selection probability vector.
Then, the following result has been established.
\begin{lemma}[Overall results in \citet{pmlr-v201-honda23a}]\label{lem: main rslt of hnd}
    In adversarial setting, FTPL with GR, Fr\'eceht distribution with shape $2$, and learning rate $\eta_t = c/\sqrt{t}$ for $c>0$ satisfies
    \begin{align*}
        \gR(T) 
        &\leq \underbrace{\sum_{t=1}^T \E\qty[\inp{\hat{\ell}_t}{w_{t+1}-w_t}]}_{\text{stability}} + \underbrace{\sum_{t=1}^T \qty(\frac{1}{\eta_{t+1}} - \frac{1}{\eta_t})\E[r_{t+1,I_{t+1}}-r_{t+1,i^*}]}_{\text{penalty}} \\
        &\hspace{9cm}+ \underbrace{\frac{\sqrt{\pi K}}{c}}_{\text{dependent term only on }\gD, \, \eta_1}\\
        &\leq \sum_{t=1}^T \E\qty[\inp{\hat{\ell}_t}{w_{t+1}-w_t}] + \frac{3.7}{c}\sqrt{KT} + \frac{\sqrt{\pi K}}{c}\\ 
        &\leq \sum_{t=1}^T \E \qty[\inp{\hat{\ell}_t}{\phi(\eta_t \hat{L}_t)-\phi(\eta_t (\hat{L}_t+\hat{\ell}_t))}] + \log(T+1)  + \frac{3.7}{c}\sqrt{KT} +\frac{\sqrt{\pi K}}{c} \numberthis{\label{eq: hnd main term}}\\
        &\leq \qty(12c\sqrt{\pi}+ \frac{3.7}{c})\sqrt{KT} + \log(T+1)+\frac{\sqrt{\pi K}}{c},
    \end{align*}
    whose dominant term can be optimized as $\gR(T)\leq 17.8\sqrt{KT} +O(\sqrt{K}+\log T)$ when $c=0.42$.
\end{lemma}
In the adversarial setting, the penalty term is bounded by $3.7\sqrt{KT}/c$, independent of the value of the cumulative loss estimator $\hat{L}_t$, as shown through worst-case analysis~\citep[see][Lemma 9]{pmlr-v201-honda23a}.
In the derivation of (\ref{eq: hnd main term}), by definition of $\phi$, we have
\begin{align*}
        w_t - w_{t+1} &= \phi(\eta_t \hat{L}_t) - \phi(\eta_{t+1} \hat{L}_{t+1}) \\
        &= \phi(\eta_t \hat{L}_t) - \phi(\eta_{t+1} (\hat{L}_{t}+\hat{\ell}_t)) \\
        &= \phi(\eta_t \hat{L}_t) - \phi(\eta_{t} (\hat{L}_{t}+\hat{\ell}_t)) + \phi(\eta_{t} (\hat{L}_{t}+\hat{\ell}_t)) - \phi(\eta_{t+1} (\hat{L}_{t}+\hat{\ell}_t)),
\end{align*}
which implies
\begin{align*}
    \sum_{t=1} \E \qty[\inp{\hat{\ell}_t}{w_t - w_{t+1}}] &=  \sum_{t=1}^T \E \qty[\inp{\hat{\ell}_t}{\phi(\eta_t \hat{L}_t)-\phi(\eta_t (\hat{L}_t+\hat{\ell}_t))}] \\
    & \qquad + \sum_{t=1}^T \E \qty[\inp{\hat{\ell}_t}{\phi(\eta_{t} (\hat{L}_{t}+\hat{\ell}_t)) - \phi(\eta_{t+1} (\hat{L}_{t}+\hat{\ell}_t))}].
\end{align*}
Here, the second term is also bounded by $\log (T+1)$, regardless of the use of GR~\citep[see][Lemma 8]{pmlr-v247-lee24a}.
Therefore, the transition from GR to CGR II affects the upper bound of the stability term, particularly when accounting for the variance of the IW estimator, which is related to the first term in (\ref{eq: hnd main term}).
The following theorem, which formalizes Theorem~\ref{thm: informal}, shows that using CGR II can further improve the dominant term.
\begin{theorem}[Formal version of Theorem~\ref{thm: informal}]\label{thm: formal}
    In adversarial setting, FTPL with CGR II, Fr\'echet distribution with shape $2$, and learning rate $\eta_t = c/\sqrt{t}$ satisfies
    \begin{multline*}
        \gR(T) \leq \qty(11.5c\sqrt{\pi}+\frac{3.7}{c}) \sqrt{KT} + \log(T+1) + \frac{\sqrt{\pi K}}{c} \\
        + \sum_{t=1}^T \sum_{i\in [K]}\E\qty[w_{t,i}\qty(1-\frac{w_{t,i}}{\sP(\gA_{t}| \hat{L}_t, I_t=i)})^{G_t}],
    \end{multline*}
    whose dominant term can be optimized as $\gR(T)\leq 17.37\sqrt{KT} +O(\sqrt{K}+\log T)$ when $c=0.43$ and $G_t \geq (K \lor 4)\log t$. 
\end{theorem}
Note that the optimized $c$ of CGR II is greater than that of GR, as shown in Lemma~\ref{lem: main rslt of hnd}, implying that FTPL with CGR II-unbaised can further improve the additive constant term.

Before the proof, we define another function introduced in \citet{pmlr-v201-honda23a} as
\begin{align}\label{eq: def of I}
    I_{i,n}(\lambda) = \int_0^\infty \frac{1}{(z+\lambda_i)^n} \exp(-\sum_{j=1}^K \frac{1}{(z+\lambda_j)^2}) \dd z>0,
\end{align}
which satisfies
\begin{equation}\label{def: phi and phi'}
    \phi_i(\lambda) = 2I_{i,3}(\ul) , \quad \phi_i'(\lambda) = -6I_{i,4}(\ul) + 4I_{i,6}(\ul).
\end{equation}
Then, the following result was established.
\begin{lemma}[Lemma 5 in \citet{pmlr-v201-honda23a}]\label{lem: hnd I ratio}
    If $\lambda_i$ is the $\sigma_i$-th smallest among $\lambda_1, \ldots, \lambda_K$ (ties are broken arbitrarily) then
    \begin{equation*}
        \frac{I_{i,4}(\ul)}{I_{i,3}(\ul)} \leq \frac{\sqrt{\pi/\sigma_i}}{2}.
    \end{equation*}
\end{lemma}
\begin{proof}[Proof of Theorem~\ref{thm: informal}]
    As discussed above, it suffices to show 
    \begin{equation*}
        \sum_{t=1}^T \E \qty[\inp{\hat{\ell}_t}{\phi(\eta_t \hat{L}_t)-\phi(\eta_t (\hat{L}_t+\hat{\ell}_t))}] \leq 11.5c\sqrt{\pi KT}.
    \end{equation*}
    Here, \citet{pmlr-v201-honda23a} showed that in their proof of Lemma 7 that
    \begin{align*}
        \phi_i(\eta_t \hat{L}_t) - \phi_i(\eta_t (\hat{L}_t+(\ell_{t,i}\widehat{w_{t,i}^{-1}})e_i)) &= \int_0^{\eta_t \ell_{t,i}\widehat{w_{t,i}^{-1}} } -\phi_i'(\eta_t \hat{L}_t + xe_i) \dd x \\
        &\leq 6 \int_0^{\eta_t \ell_{t,i}\widehat{w_{t,i}^{-1}} } I_{i,4}(\underline{\eta_t \hat{L}_t + xe_i}) \dd x \tag{by (\ref{def: phi and phi'})} \\
        &\leq 6 \int_0^{\eta_t \ell_{t,i}\widehat{w_{t,i}^{-1}} } I_{i,4}(\eta_t \hat{\uL}_t) \dd x \tag{monotonicity of $I_{i,4}$} \\
        &\leq 6 \eta_t \ell_{t,i}  I_{i,4}(\eta_t \hat{\uL}_t)  \widehat{w_{t,i}^{-1}}.
    \end{align*}
    Here, the monotonicity of $I_{i,4}$ in its $i$-th element might not appear trivial.
    From the definition of $I$ in (\ref{eq: def of I}), one can interpret this function as an arm-selection probability of FTPL when $r_i$ is of density of order $(z+\lambda_i)^4$ and $r_j$s for $j\ne i$ follows the Fr\'echet distribution with shape $2$.
    Therefore, $I_{i,4}$ is monotonically decreasing with respect to the $i$-th argument.
    
    Under CGR II, we obtain
    \begin{align*}
        \E\qty[\widehat{w_{t,I_t}^{-1}}^2 \middle | \hat{L}_t, I_t] &= \mathrm{Var}\qty[\widehat{w_{t,I_t}^{-1}}\middle | \hat{L}_t, I_t] + \E^2\qty[\widehat{w_{t,I_t}^{-1}} \middle | \hat{L}_t, I_t] \\
        &= \frac{2}{w_{t,I_t}^2} - \frac{1}{w_{t,I_t}\sP(A_{t})}.
    \end{align*}
    Recall (\ref{eq: P/w ratio}), which shows
    \begin{equation*}
        \frac{\sP(A_{t,i})}{w_{t,i}} \leq \frac{4h(x)}{h(y)},
    \end{equation*}
    where $h(x) = \frac{1-e^{-x}}{x}$, $x=\frac{\sigma_i}{\ul_{t,i}^2}$, and $y=\frac{3(\sigma_i-1)+K}{4\ul_{t,i}^2}$.
    One can see that for any $i \in [K]$ and $K\geq 2$
    \begin{equation*}
        \frac{y}{x} = \frac{3(\sigma_i -1) +K}{4\sigma_i} \leq \frac{K}{\sigma_i}.
    \end{equation*}
    Following the same arguments in (\ref{eq: h ratio}), we have
    \begin{equation*}
       \frac{\sP(A_{t,i})}{w_{t,i}} \leq \frac{4h(x)}{h(y)} \leq \sup_{x>0} \frac{4h(x)}{h(Kx/\sigma_i)} \leq \frac{4K}{\sigma_i},
    \end{equation*}
    which implies 
    \begin{equation*} % \label{eq: ub of 2nd moment}
         \frac{1}{\sP(A_t)} \geq \frac{\sigma_i}{4Kw_{t,i}} \implies \E\qty[\widehat{w_{t,I_t}^{-1}}^2 \middle | \hat{L}_t, I_t] \leq \qty(2-\frac{\sigma_i}{4K})\frac{1}{w_{t,I_t}^2}.
    \end{equation*}
    Since $\hat{\ell}_t = (\ell_{t,i}\widehat{w_{t,i}^{-1}})e_{i}$ when $I_t = i$, we obtain
    \begin{align*}
       \E \qty[\hat{\ell}_{t,i}( \phi_i(\eta_t \hat{L}_t) - \phi_i(\eta_t (\hat{L}_t + \hat{\ell}_t)) )\middle | \hat{L}_t] &\leq \E\qty[\I[I_t =i]\ell_{t,i}\widehat{w_{t,i}^{-1}} \cdot 6\eta_t \ell_{t,i}  I_{i,4}(\eta_t \hat{\uL}_t)  \widehat{w_{t,i}^{-1}} \middle | \hat{L}_t] \\
       &\leq 6\eta_t \E\qty[\qty(2-\frac{\sigma_i}{4K})w_{t,i} \frac{\ell_{t,i}^2 I_{i,4}(\eta_t \hat{\uL}_{t,i})}{w_{t,i}^2} \middle | \hat{L}_t] \\
       &= 3\eta_t \E\qty[\qty(2-\frac{\sigma_i}{4K})\frac{\ell_{t,i}^2 I_{i,4}(\eta_t \hat{\uL}_{t,i})}{I_{i,3}(\eta_t \hat{\uL}_{t,i})} \middle | \hat{L}_t] \tag{by (\ref{def: phi and phi'})} \\
       &\leq 3\eta_t \E\qty[\qty(2-\frac{\sigma_i}{4K})\frac{I_{i,4}(\eta_t \hat{\uL}_{t,i})}{I_{i,3}(\eta_t \hat{\uL}_{t,i})} \middle | \hat{L}_t] \tag{$\because \ell_{t,i}\in [0,1]$} \\
       &\leq \frac{3\sqrt{\pi}}{2}\eta_t \E\qty[\qty(\frac{2}{\sqrt{\sigma_i}}-\frac{\sqrt{\sigma_i}}{4K}) \middle | \hat{L}_t]. \tag{by Lemma~\ref{lem: hnd I ratio}}
    \end{align*}
    Therefore,
    \begin{align*}
       \E \qty[\inp{\hat{\ell}_t}{\phi(\eta_t \hat{L}_t)-\phi(\eta_t (\hat{L}_t+\hat{\ell}_t))}] &\leq  \frac{3\sqrt{\pi}}{2}\eta_t\sum_{i=1}^K \qty(\frac{2}{\sqrt{\sigma_i}}- \frac{\sqrt{\sigma_i}}{4K}) \\
       &\leq 3\sqrt{\pi}\eta_t \qty(1+\int_1^K x^{-1/2} \dd x)  - \frac{3\sqrt{\pi}}{2}\eta_t \qty(\int_0^K \frac{x^{1/2}}{4K} \dd x) \\
       &= 3\sqrt{\pi}\eta_t(2\sqrt{K}-1) - \sqrt{\pi}\eta_t \frac{\sqrt{K}}{4} 
    %    \numberthis{\label{eq: improvement}} 
       \\
       &\leq \frac{23\sqrt{\pi}}{4}\sqrt{K}\eta_t.
    \end{align*}
    By taking summation with $\eta_t = c/\sqrt{t}$, we have
    \begin{align*}
         \sum_t \E \qty[\inp{\hat{\ell}_t}{\phi(\eta_t \hat{L}_t)-\phi(\eta_t (\hat{L}_t+\hat{\ell}_t))}] &\leq \frac{23\sqrt{\pi}c}{4}\sqrt{K} \sum_{t=1}\frac{1}{\sqrt{t}} \\
         &\leq  \frac{23\sqrt{\pi}c}{2}\sqrt{KT},
    \end{align*}
    which concludes the proof.
\end{proof}
